\documentclass[11pt]{article}
\usepackage[margin=1.15in]{geometry}
\usepackage{amsmath,amssymb,amsthm}
\usepackage{hyperref}
\newtheorem{theorem}{Theorem}
\newtheorem{proposition}[theorem]{Proposition}
\newtheorem{problem}[theorem]{Problem}
\theoremstyle{remark}
\newtheorem{remark}[theorem]{Remark}

\title{A graph with four card types and five vertex orbits:\\
the $k=4$ case of Kourovka Notebook Problem 20.71}
\author{Jitendra Prajapati\thanks{IIT Madras. Verification scripts and a
Lean~4 certificate: \url{https://github.com/infinityscroll/kourovka-20-71}.}}
\date{July 2026 (draft)}

\begin{document}
\maketitle

\begin{abstract}
A \emph{card} of a finite simple graph $G$ is a vertex-deleted induced
subgraph $G - v$. Problem 20.71 of the Kourovka Notebook (J.~Lauri,
communicated by V.~Pannone) asks, for $k = 2, 3, 4$, whether a connected
graph with exactly $k$ isomorphism types of cards can have more than $k$
orbits of $\operatorname{Aut}(G)$ on vertices. We answer the case $k = 4$
affirmatively with a graph on $9$ vertices, which an exhaustive census shows
to be of minimum order and unique up to complementation; curiously, no
example of order $10$ exists for any $k \in \{2,3,4\}$. The verification is
elementary, and is additionally certified in Lean~4. The same census shows that any example for $k = 2$ or
$k = 3$ must have at least $12$ vertices.
\end{abstract}

\section{Introduction}

Let $G$ be a finite simple undirected graph with vertex set $V(G)$,
$|V(G)| = n$. The \emph{cards} of $G$ are the $n$ induced subgraphs
$G - v$, $v \in V(G)$. Vertices in one orbit of $\operatorname{Aut}(G)$
yield isomorphic cards, so the partition of $V(G)$ by card type is coarser
than the orbit partition; in particular a graph with exactly $k$ card types
has at least $k$ orbits, and vertices with isomorphic cards lying in
distinct orbits are the classical \emph{pseudosimilar} vertices
\cite{GK82,Lauri97,LS16}. If all cards are mutually isomorphic ($k=1$), it
is known that $G$ must be vertex-transitive, i.e.\ the orbit count is
exactly $1$ \cite[Problem 20.71, comment]{KN21}.

Problem 20.71 of the Kourovka Notebook \cite{KN21}, due to J.~Lauri and
communicated by V.~Pannone, asks:

\begin{problem}[{\cite[20.71]{KN21}}]\label{prob}
For $k = 2, 3, 4$: can a connected graph with exactly $k$ isomorphism types
of cards have more than $k$ orbits of $\operatorname{Aut}(G)$ on $V(G)$?
\end{problem}

Examples are recorded there with $5$ card types and $6$ orbits, and with
$6$ card types and $7$ orbits. We settle the case $k = 4$:

\begin{theorem}\label{thm:main}
There is a connected graph $W$ on $9$ vertices with exactly four
isomorphism types of cards and exactly five vertex orbits.
\end{theorem}

\begin{proposition}[computational]\label{prop:census}
By exhaustive search over connected unlabeled graphs:
\begin{enumerate}
\item no graph of order at most $8$ has $k \in \{2,3,4\}$ card types and
more than $k$ orbits, so $W$ has minimum order;
\item at order $9$ there are exactly two such graphs for $k = 4$, namely
$W$ and its complement $\overline{W}$ (card types and orbits are invariant
under complementation, and $\overline{W}$ is again connected), so the
example is unique up to complementation; at order $10$ there is none for
any $k \in \{2,3,4\}$;
\item no graph of order at most $11$ has $k \in \{2,3\}$ card types and
more than $k$ orbits. Consequently any affirmative example for the two
remaining cases of Problem~\ref{prob} has at least $12$ vertices.
\end{enumerate}
\end{proposition}

\section{The witness}

Let $W$ be the graph on $\{0, 1, \dots, 8\}$ with the $16$ edges
\[
03,\ 04,\ 07,\ 14,\ 15,\ 16,\ 25,\ 26,\ 27,\ 28,\ 36,\ 37,\ 38,\ 48,\ 57,\ 68
\]
(graph6 code \verb|HCpbdgy|). $W$ is connected: $03, 04, 07, 14, 15, 16,
25, 28$ is a spanning tree.

\subsection{Exactly five orbits}

Colour refinement starting from degrees stabilises at the partition
\[
\pi \;=\; \{0,5\} \mid \{1,4\} \mid \{2,3\} \mid \{6,8\} \mid \{7\},
\]
and every automorphism preserves the cells of $\pi$. The involution
\[
\sigma = (0\,5)(1\,4)(2\,3)(6\,8)
\]
is an automorphism of $W$ (a direct check on the edge list), so each cell
of $\pi$ is a single orbit: there are exactly five orbits, and
$\operatorname{Aut}(W) = \{\mathrm{id}, \sigma\}$ (the exhaustive check
over all $9!$ permutations confirms no further automorphisms, though for
Theorem~\ref{thm:main} the count of orbits needs only $\pi$ and $\sigma$).

\subsection{Exactly four card types}

The degree multisets of the nine cards take exactly four values:
\[
\begin{array}{ll}
W - 0,\ W - 5: & (2,3,3,3,3,4,4,4) \\
W - 1,\ W - 4: & (2,2,3,3,4,4,4,4) \\
W - 2,\ W - 3,\ W - 6,\ W - 8: & (2,3,3,3,3,3,3,4) \\
W - 7: & (2,2,3,3,3,3,4,4)
\end{array}
\]
so there are at least four card types. Conversely $\sigma$ identifies the
cards within the pairs $\{0,5\}$, $\{1,4\}$, $\{2,3\}$, $\{6,8\}$, and the
explicit bijections
\[
\begin{array}{ll}
W-2 \to W-6: & 0\mapsto2,\ 1\mapsto4,\ 3\mapsto7,\ 4\mapsto8,\ 5\mapsto1,\
6\mapsto0,\ 7\mapsto5,\ 8\mapsto3,\\
W-2 \to W-8: & 0\mapsto3,\ 1\mapsto1,\ 3\mapsto7,\ 4\mapsto6,\ 5\mapsto4,\
6\mapsto5,\ 7\mapsto0,\ 8\mapsto2,
\end{array}
\]
are isomorphisms (checks on the edge lists), so the four vertices
$2, 3, 6, 8$ give one card type. Hence there are exactly four card types,
with deletion classes $\{0,5\}, \{1,4\}, \{2,3,6,8\}, \{7\}$.

The point of the example is that the card class $\{2,3,6,8\}$ is the union
of the \emph{two} orbits $\{2,3\}$ and $\{6,8\}$: these four vertices are
pairwise pseudosimilar across the two orbits. This proves
Theorem~\ref{thm:main}. \qed

\begin{remark}
The verification above is deliberately software-independent: one involution,
two explicit card isomorphisms, four degree multisets, and colour
refinement, all checkable by hand. A machine-checked Lean~4 certificate of
the statement (connectivity, card-type count $4$, orbit count $5$), which
compiles against Mathlib with no unproven assumptions, is provided in the
repository, together with scripts reproducing every claim.
\end{remark}

\section{The census}\label{sec:census}

Proposition~\ref{prop:census} rests on exhaustive generation of connected
unlabeled graphs with \textsf{nauty}'s \texttt{geng} \cite{MP14}, in the
counts recorded at OEIS A001349 ($11{,}716{,}571$ graphs of order $10$;
$1{,}006{,}700{,}565$ of order $11$), filtered in two sound pruning stages:
a graph with at most $3$ card types has at most $3$ distinct degrees
(isomorphic cards force equal deleted-vertex degrees), and the partition of
cards by the invariant pair (degree multiset, triangle count) is coarser
than the partition into isomorphism types. Cards of the graphs surviving
both filters were classified exactly by canonical labelling
(\textsf{pynauty}), and orbit counts of the resulting candidates computed
exactly. At orders $10$ and $11$ respectively, $741{,}839$ and
$27{,}219{,}751$ graphs passed the degree filter, $8{,}059$ and $60{,}536$
required canonical labelling, and none was an example for $k \in \{2,3\}$.
Structured sweeps beyond order $11$ also found no example: all regular
graphs of orders $12$--$14$ (via complement invariance it suffices to
generate degrees $d \le (n-1)/2$; the largest slice comprises the
$21{,}609{,}300$ connected $6$-regular graphs of order $14$), and all graphs
of orders $12$--$13$ whose degrees lie in a window $\{a, a+1\}$,
$2 \le a \le 5$ ($15{,}409{,}504$ and $836{,}922{,}152$ graphs
respectively; complementation extends the coverage to the mirrored windows).
Full data are in the repository.

\section{Remaining cases}

The cases $k = 2$ and $k = 3$ of Problem~\ref{prob} remain open; by
Proposition~\ref{prop:census} any example has order at least $12$. It seems
plausible that the difficulty is genuine: fewer card types force the graph
closer to vertex-transitivity, while extra orbits require asymmetry, and
the tension is sharpest at $k = 2$ where the $k = 1$ rigidity theorem sits
immediately below.

\section*{Disclosure}
The search, the write-up and the Lean formalization were prepared with
substantial assistance from Claude (Anthropic); the witness was
additionally re-verified by independent brute-force computation, and the
central claim is machine-checked.

\begin{thebibliography}{9}
\bibitem{GK82} C.~D.~Godsil, W.~L.~Kocay, \emph{Constructing graphs with
pairs of pseudo-similar vertices}, J.~Combin.\ Theory Ser.~B 32 (1982).
\bibitem{KN21} E.~I.~Khukhro, V.~D.~Mazurov (eds.), \emph{Unsolved Problems
in Group Theory: The Kourovka Notebook}, No.~21, Sobolev Institute of
Mathematics, Novosibirsk, 2026.
\bibitem{Lauri97} J.~Lauri, \emph{Pseudosimilarity in graphs --- a survey},
Ars Combin.\ 46 (1997).
\bibitem{LS16} J.~Lauri, R.~Scapellato, \emph{Topics in Graph Automorphisms
and Reconstruction}, 2nd ed., LMS Lecture Note Series 432, Cambridge
University Press, 2016.
\bibitem{MP14} B.~D.~McKay, A.~Piperno, \emph{Practical graph isomorphism,
II}, J.~Symbolic Comput.\ 60 (2014) 94--112.
\end{thebibliography}

\end{document}
